% 原书物理页=印刷页33--35；1.10 全节.（本节自原书第 33 页下部的 1.10 节标题起；第 33 页上部属 1.9 节，已由 ch01-s09.tex 译出，此处不再重复.）
\section{不变测度及其应用}\label{sec:lie-1-10}

\emph{回顾.}\ 局部紧空间 $X$ 上的一个 \newterm{lie-radon-measure}{Radon 测度}{Radon measure} 是 $\mathcal C_c(X)$（具有紧支撑的连续函数）上的一个线性形式 $\mu$，它在下述意义下是正的：$f\geqslant 0\Rightarrow\mu(f)\geqslant 0$. 我们不加区别地记
\[
\mu(f)=\int_X f\,d\mu=\int_X f(x)\,d\mu(x).
\]
如果一个群 $G$ 作用在 $X$ 上，就让它按公式 $gf(x):=f(g^{-1}x)$ 作用在 $\mathcal C_c(X)$ 上，并按公式 $g\mu(f)=\mu(g^{-1}f)$ 作用在\newterm{lie-distribution-measure}{分布}{distribution}上，于是我们有
\[
\int_X f(x)\,d(g\mu)(x)=\int_X(g^{-1}f)(x)\,d\mu(x)=\int_X f(gx)\,d\mu(x).
\]
如果对一切 $g$ 都有 $g\mu=\mu$，就称分布 $\mu$ 是\newterm{lie-invariant-measure}{不变的}{invariant}.

\begin{theorem}\label{thm:lie-1-10-1}
在一个 Lie 群 $G$ 上，存在左不变且非零的 Radon 测度 $\mu$，并且它在相差一个正因子的意义下是唯一的.
\end{theorem}

\begin{remark*}
这个结果对任何局部紧群都成立，并且以 Haar 定理之名为人所知. 测度 $\mu$ 称为（左）\newterm{lie-haar-measure}{Haar 测度}{Haar measure}.
\end{remark*}

\begin{proof}
记 $n$ 为 $G$ 的维数. 取 $\operatorname{Lie}(G)$ 上一个非零的交错 $n$-形式 $dX$. 它也是 $\operatorname{Lie}(G)$ 上的一个 Lebesgue 测度（取定合适的定向）. 对每个 $g\in G$，令 $\omega(g):=L_g^{*}(dX)$，它是 $T_gG$ 上的一个交错 $n$-形式. 这样我们便得到 $G$ 上的一个 \newterm{lie-volume-form}{体积形式}{volume form} $\omega$，即典范丛 $\Omega G=\bigwedge^{n}T^{*}G$ 的一个处处不为零的光滑截面，并且它在左平移下不变：对一切 $g\in G$ 有 $L_g^{*}\omega=\omega$. 流形上的积分理论于是把一个正测度 $\mu$ 与 $\omega$ 联系起来，它与 $\omega$ 一样在左平移下不变.
在我们的情形，可以按如下方式构造 $\mu$. 设 $\mathcal U$ 是 $\operatorname{Lie}(G)$ 中 $0$ 的一个开邻域，使得 $\exp_G$ 在其上实现一个微分同胚 $\mathcal U\xrightarrow{\ \sim\ }\mathcal V=\exp(\mathcal U)$. $\mathcal U$ 上的测度 $\exp^{*}(\omega|_{\mathcal V})$ 具有形状 $\varepsilon(X)\,dX$，其中 $\varepsilon$ 是 $\mathcal U$ 上一个处处不为零、在 $0$ 处取值为 $1$ 的光滑函数，因而是（严格）正的. 设 $f\in\mathcal C_c(G)$，其支撑 $\operatorname{Supp}(f)$ 包含在 $\mathcal V$ 之中. 此时令
\[
\mu(f):=\int_{\mathcal U}f(\exp(X))\exp^{*}(\omega)=\int_{\mathcal U}f(\exp(X))\varepsilon(X)\,dX.
\]
如果 $g\in G$ 满足 $\operatorname{Supp}(gf)=g\operatorname{Supp}(f)\subset\mathcal V$，那么 $\omega$ 在 $L_g$ 下的不变性蕴含 $\mu(f)=\mu(gf)$. 因此，更一般地，若从某个 $f\in\mathcal C_c(G)$ 出发，其支撑包含在某个 $g\mathcal V$（其中 $g\in G$）之中，则可以令 $\mu(f):=\mu(g^{-1}f)$，而这不依赖于 $g$ 的选取. 最后，对任意的 $f\in\mathcal C_c(G)$，可以利用 \newterm{lie-partition-of-unity}{单位分解}{partition of unity} 把它写成 $f=\sum_{i=1}^{m}f_i$ 的形式，其中每个 $\operatorname{Supp}(f_i)$ 都包含在 $\mathcal V$ 的某个平移之中. 那么，表达式 $\mu(f):=\sum_{i=1}^{m}\mu(f_i)$ 不依赖于这种写法，从而定义了所要求的左不变 Radon 测度 $\mu$.
还可以把 $\exp^{*}(\omega)$ 在 $\mathcal U$ 上、或至少在一个对称开集 $\mathcal U'\subset\mathcal U$（满足 $(\exp(\mathcal U'))^{2}\subset\mathcal V$）上的形状写得更精确. 事实上，对 $X\in\mathcal U'$，考虑映射 $\alpha_X:\mathcal U'\longrightarrow\mathcal U$，$Y\mapsto\exp^{-1}(\exp(X)\exp(Y))$. 这是 $\mathcal U'$ 到它的像上的一个微分同胚，它把 $0$ 送到 $X$. 按照 $\omega$ 的定义进行计算，可求得 $\exp^{*}(\omega)=\det(d_0(\alpha_X))\,dX$.
这样我们就证明了存在性. 至于唯一性，可以注意到 $\Omega_G=\bigwedge^{n}(T^{*}G)$ 的不变截面所构成的空间显然是 $1$ 维的.
\end{proof}

\noindent\textbf{例.}\ 记 $dx$ 为 $\mathbb R$ 上或有限维 $\mathbb R$-向量空间上的 Lebesgue 测度.
\begin{enumerate}[label=\roman*)]
\item 在 $G=\mathbb R$ 上，测度 $dx$ 是不变的（左不变且右不变）.
\item 在 $G=\mathbb R^{\times}$ 上，测度 $d^{\times}x:=\dfrac{dx}{|x|}$ 是不变的.
\item 在 $G=S^{1}$ 上，``该''不变测度由 $\int_{S^{1}}f\,\mu=\int_{0}^{1}f(\exp(2i\pi\theta))\,d\theta$ 给出.
\item 在 $G=\mathbb C^{\times}$ 上，测度 $\dfrac{i}{2}\dfrac{dz\wedge d\bar z}{z\bar z}=\dfrac{dx\,dy}{x^{2}+y^{2}}$ 是不变的.
\item 在 $G=\mathrm{GL}_n(\mathbb R)^{0}$ 上，测度 $\dfrac{dx}{|\det(x)|^{n}}$ 是左不变且右不变的. 事实上，我们先让 $G$ 作用在 $\mathbb R^{n}$ 上，并记 $dx=dx_1\cdots dx_n$ 为 $\mathbb R^{n}$ 上的 Lebesgue 测度. 那么，众所周知
\[
\int_{\mathbb R^{n}}f(gx)\,dx=\int_{\mathbb R^{n}}f(x)|\det(g)|^{-1}\,dx,\quad\forall f\in\mathcal C_c(\mathbb R^{n}).
\]
现在让 $G$ 通过左乘作用在 $\mathrm{M}_n(\mathbb R)$ 上，并记 $dx=dx_{11}\cdots dx_{n1}dx_{12}\cdots\cdots dx_{nn}$ 为 $\mathrm{M}_n(\mathbb R)$ 上的 Lebesgue 测度. 那么，把 $\mathrm{M}_n(\mathbb R)=\mathbb R^{n}\oplus\cdots\oplus\mathbb R^{n}$ 按列分解，便得到
\[
\int_{\mathrm{M}_n(\mathbb R)}f(gx)\,dx=\int_{\mathrm{M}_n(\mathbb R)}f(x)|\det(g)|^{-n}\,dx,\quad\forall f\in\mathcal C_c(\mathrm{M}_n(\mathbb R)).
\]
由此可得\footnote{译注：下式中第一个等号右端积分号后的表达式按原文照录；按计算此处应为 $|\det(g^{-1}x)|^{-n}\,|\det(g)|^{-n}$，原文排印时绝对值符号的位置有误.}
\begin{align*}
\int_{G}f(gx)|\det(x)^{-n}|\,dx
&=\int_{G}f(x)\det(g^{-1}x)^{-n}||\det(g)|^{-n}|\,dx\\
&=\int_{G}f(x)|\det(x)^{-n}|\,dx,\quad\forall f\in\mathcal C_c(G).
\end{align*}
\end{enumerate}

\begin{remark*}
如果 $\varphi:G'\xrightarrow{\ \sim\ }G$ 是一个同构，而 $\omega$ 是 $G$ 上一个左不变的体积形式，那么 $\varphi^{*}(\omega)$ 显然是 $G'$ 上一个左不变的体积形式. 这适用于 $G'=G$ 与 $\varphi=\operatorname{Int}_g$ 的情形. 由左不变体积形式在相差一个纯量意义下的唯一性，得到 $\operatorname{Int}_g^{*}(\omega)=\lambda(g)\omega$，且映射 $g\mapsto\lambda(g)\in\mathbb R^{\times}$ 是乘性的. 为计算 $\lambda(g)$，我们在点 $e$ 处考察. 我们有 $\operatorname{Int}_g^{*}(\omega)(e)=\operatorname{Ad}_g^{*}(\omega(e))=\operatorname{Ad}_g^{*}(dX)=\det(\operatorname{Ad}_g)\,.\,dX$，故 $\lambda(g)=\det(\operatorname{Ad}_g)$. 过渡到不变测度上，便得到 $\operatorname{Int}_g^{*}(\mu)=|\det(\operatorname{Ad}_g)|\,\mu$.
\end{remark*}

\begin{numpara}[右不变测度]\label{np:lie-1-10-2}
上面的定理对右不变测度显然依然成立. 但前面的注给出了一个从左过渡到右（！）的办法. 事实上，若 $\omega$ 是一个左不变的体积形式，则有
\[
R_g^{*}(\omega)=R_g^{*}L_g^{*}(\omega)=\operatorname{Int}_g^{*}(\omega)=\det(\operatorname{Ad}_g)\,\omega.
\]
于是我们得到了：
\begin{proposition*}
若 $\mu$ 是一个左不变测度，则测度 $|\det(\operatorname{Ad})|\,.\,\mu$ 是右不变的.
\end{proposition*}
如果 $|\det(\operatorname{Ad})|\equiv 1$，就称 $G$ 是 \newterm{lie-unimodular}{幺模的}{unimodular}. 可以证明各个经典群都是幺模的，而最简单的例子是下面这个：
\begin{example*}
设 $G$ 紧. 由于 $G$ 在 $|\det\operatorname{Ad}|$ 下的像是 $\mathbb R^{\times}_{>0}$ 的一个紧子群，故 $|\det(\operatorname{Ad})|\equiv 1$. 因而紧群 $G$ 上的任何左不变测度 $\mu$ 都是右不变的. 通常把 $\mu$ 规范化为 $\mu(G)=\mu(1_G)=1$.
\end{example*}
\end{numpara}

\begin{numpara}[应用]\label{np:lie-1-10-3}
不变测度之存在性的许多应用都是下述论断的各种变形：
\begin{lemma*}
设 $G$ 紧，并记 $\mu$ 为 $G$ 上规范化后的 Haar 测度. 设 $\rho:G\longrightarrow\mathrm{GL}(V)$ 是一个表示，$v\in V$. 那么 $\bar v:=\int_{G}\rho(x)(v)\,d\mu(x)$ 是 $V$ 中一个 $G$-不变向量.
\end{lemma*}
\begin{proof}
事实上，对一切 $g\in G$，有 $\rho(g)(\bar v)=\int_{G}\rho(gx)(v)\,d\mu(x)=\int_{G}\rho(x)(v)\,d\mu(x)=\bar v$.
\end{proof}
\begin{example*}
设 $G$ 是 $\mathrm{GL}(V)$ 的一个紧子群，$\langle\ ,\ \rangle$ 是 $V$ 上一个欧几里得内积. 令
\[
\Phi(v,w):=\int_{G}\langle xv,xw\rangle\,d\mu(x).
\]
这是一个对称双线性型，并且是正定的，因为 $\langle\ ,\ \rangle$ 是正定的. 此外，它在 $G$ 下不变，这也意味着 $G\subset\mathrm{O}(\Phi)$. 由此可得：
\end{example*}
\begin{corollary*}
$\mathrm{GL}_n(\mathbb R)$ 的任何紧子群都共轭于 $\mathrm{O}(n,\mathbb R)$ 的一个子群.
\end{corollary*}
\begin{exercise*}
证明 $\mathrm{GL}_n(\mathbb C)$ 的任何紧子群都共轭于 $\mathrm{U}(n)$ 的一个子群.
\end{exercise*}
\end{numpara}
